\section*{Bab \ref{ch:b2:deep-time-climate} --- Iklim sepanjang Waktu Geologi}

\begin{solution}{exo:b2:deep-time-climate:1}
(a) $900 + 37100 + 700 + 3 = \qty{38703}{PgC}$, sekitar 65 kali \qty{591}{PgC} di udara.
(b) $\num{3.5e6}\times0.012011 = \qty{42000}{PgC}$. (c) $591/42000 = \qty{1.4}{\percent}$.
\end{solution}

\begin{solution}{exo:b2:deep-time-climate:2}
(a) Masukan $11.3 + 3 + 12.3 + 8.1 + 2.0 = \qty{36.7}{Tmol.yr^{-1}}$; keluaran $26 + 16 =
\qty{42}{Tmol.yr^{-1}}$, sama di dalam batas ketidakpastiannya. (b) $(11.3 + 3)/36.7 = 0.39$: jauh
di bawah separuh; sebagian besar masukan adalah karbon sedimen tua. (c) $\num{3.5e6}/42 = \qty{8.3e4}{yr}$.
\end{solution}

\begin{solution}{exo:b2:deep-time-climate:3}
(a) \ce{Mg2SiO4 + 4 CO2 + 2 H2O -> 2 Mg^{2+} + 4 HCO3^- + SiO2}. (b) Dengan \ce{2 Mg^{2+} + 4 HCO3^-
-> 2 MgCO3 + 2 CO2 + 2 H2O}, reaksi bersihnya adalah \ce{Mg2SiO4 + 2 CO2 -> 2 MgCO3 + SiO2}: dua mol
karbon dioksida per mol forsterit. (c) $\qty{1}{PgC} = 10^{15}/12.011 =
\qty{8.33e13}{mol}$ karbon, yang memerlukan \qty{4.16e13}{mol} forsterit, $4.16\times10^{13}
\times\qty{140.7}{g} = \qty{5.86}{Pg}$, hampir enam miliar ton.
\end{solution}

\begin{solution}{exo:b2:deep-time-climate:4}
Glasiasi Hirnantium, 445 sampai \qty{443}{Ma}, rumah es; luas es terbesar pada Paleozoikum akhir, 307 sampai
\qty{299}{Ma}, rumah es; Turonium, 94 sampai \qty{89}{Ma}, rumah kaca (masa terhangat dalam Fanerozoikum);
PETM, 55.9 sampai \qty{55.7}{Ma}, rumah kaca; EECO, 53 sampai \qty{49}{Ma}, rumah kaca; awal lembaran
es Antarktika, sekitar \qty{34}{Ma}, permulaan rumah es masa kini; lembaran es utara,
\qty{2.7}{Ma}, rumah es.
\end{solution}

\begin{solution}{exo:b2:deep-time-climate:5}
(a) $63000/(8.314\times288^2) = \qty{0.0914}{K^{-1}}$, yang dibulatkan menjadi 0.09. (b) $74000/(8.314
\times288^2) = \qty{0.107}{K^{-1}}$. (c) $\ln1.2/0.09 = \qty{2.03}{K}$ dan $\ln 1.2/0.107 =
\qty{1.70}{K}$. (d) Pada keadaan tunak pelapukan harus sama dengan pelepasan gas, yang menetapkan $e^{k
\Delta T}$ dan dengan demikian $\Delta T$; sensitivitas hanya menentukan berapa banyak karbon dioksida yang menghasilkan
pemanasan itu.
\end{solution}

\begin{solution}{exo:b2:deep-time-climate:6}
Di sini $k_T = 0.0914$. (a) Eksak $\exp[7577.5\times(1/288 - 1/293)] = 1.567$; pendekatan
$e^{0.0914\times5} = 1.579$: \qty{0.8}{\percent} terlalu tinggi. (b) Eksak $e^{2.482} = 11.97$;
pendekatan $e^{2.741} = 15.50$: \qty{30}{\percent} terlalu tinggi. (c) $\ln(\text{rasio}) = (E_{\text
a}/R)[x/T_0^2 - x/(T_0(T_0 + x))] = (E_{\text a}/R)\,x^2/(T_0^2(T_0 + x))$. Dengan menyamakannya dengan $\ln1.1$:
$x^2 = 1.043\,(288 + x)$, sehingga $x = \qty{17.9}{K}$. Di atas pemanasan sekitar \qty{18}{K}
bentuk eksponensial itu melebih-lebihkan pelapukan dengan lebih dari sepersepuluh.
\end{solution}

\begin{solution}{exo:b2:deep-time-climate:7}
(a) $1444/276 = 5.23$, dan $3.93\log_2 5.23 = \qty{9.4}{W.m^{-2}}$. (b) Hampir tepat sama dengan
\qty{9}{W.m^{-2}} dari Matahari: kedua gaya radiatif itu saling meniadakan. (c) $(1444 - 276)/419.5 =
\qty{2.8}{ppm}$ per juta tahun. Nilai \qty{3.4}{ppm} berasal dari garis lurus yang dicocokkan dengan semua
data, yang memberi bobot pada setiap perkiraan, bukan pada kedua ujung kurva yang dihaluskan.
\end{solution}

\begin{solution}{exo:b2:deep-time-climate:8}
(a) Dari $3000/20000 = \qty{0.15}{PgC.yr^{-1}}$ sampai $7000/3000 = \qty{2.3}{PgC.yr^{-1}}$, rentang
yang mencakup 0.3 sampai 1.5 dari IPCC. (b) $35\times12.011/44.01 = \qty{9.6}{PgC.yr^{-1}}$, 6 sampai
32 kali 0.3 sampai \qty{1.5}{PgC.yr^{-1}} pada PETM. (c) $3.93\log_2(2000/900) = \qty{4.5}{W.m^{-2}}$,
dan $3\log_2 2.22 = \qty{3.5}{\celsius}$: sekitar tujuh puluh persen dari \qty{5}{\celsius}. Umpan balik
lain atau sumber pemanasan lain, seperti metana, harus menutupi sisanya.
\end{solution}

\begin{solution}{exo:b2:deep-time-climate:9}
(a) Sekarang $1.2\,W_0e^{k\Delta T} = V = W_0$, sehingga $\Delta T = -\ln 1.2/0.09 = \qty{-2.03}{K}$.
(b) $278\times2^{-2.03/3} = \qty{174}{ppm}$. (c) Tidak: pada keadaan tunak mana pun pelapukan sama dengan
pelepasan gas, yang tidak berubah. Kemampuan lapuk yang lebih tinggi diimbangi oleh iklim yang lebih dingin.
\end{solution}

\begin{solution}{exo:b2:deep-time-climate:10}
(a) $k = 0.119$: $\tau = 591\times0.693/(0.08\times0.119\times3\times0.138) = \qty{104}{ka}$.
(b) $S = 8$: $\tau = \qty{51.5}{ka}$. (c) Keduanya membuat termostat lebih cepat, karena $\tau \propto
1/(kS)$. $E_{\text a}$ tidak pasti dengan faktor sekitar dua, dan $S$ pada skala waktu ini dengan faktor
yang serupa (3 dibandingkan dengan 5 sampai \qty{8}{K}): keduanya kira-kira sama pentingnya.
\end{solution}

\begin{solution}{exo:b2:deep-time-climate:11}
(a) $0.08\times5000 = \qty{400}{PgC}$ tetap berada di udara: $C/C_0 = 991/591 = 1.677$, atau
\qty{466}{ppm}; $\Delta T = 3\log_2 1.677 = \qty{2.24}{K}$. (b) $W - W_0 = 0.138\,(e^{0.09\times2.24}
- 1) = \qty{0.0308}{PgC.yr^{-1}}$, pelapukan \qty{22}{\percent} lebih banyak; $5000/0.0308 =
\qty{162}{ka}$. (c) $\tau = \qty{137}{ka}$ dan $\tau\ln10 = \qty{315}{ka}$. (d) Pemanasan tumbuh
hanya menurut logaritma karbon dioksida, sehingga kelebihan yang besar menaikkan pelapukan lebih sedikit daripada
yang diandaikan hukum linear; pemindahannya mula-mula lebih lambat.
\end{solution}

\begin{solution}{exo:b2:deep-time-climate:12}
(a) Di dekat $M^*$, $\mathrm d\Delta T/\mathrm dM = Sf/(A_0\ln 2\,(1 + fM^*/A_0)) = Sf/(A_0\ln 2\,
C^*/C_0)$ dan $\mathrm dW/\mathrm dM = kW\,\mathrm d\Delta T/\mathrm dM$, dengan $W = V$ di sana;
kebalikannya adalah $\tau'$. (b) $C^*/C_0 = 2^{2.03/3} = 1.597$: $\tau' = 137.4\times1.597/1.2 =
\qty{183}{ka}$. (c) Lebih lambat: pada karbon dioksida yang tinggi setiap petagram karbon tambahan menghangatkan lebih sedikit,
karena gaya radiatifnya logaritmik, sehingga pelapukan lebih sedikit menanggapinya.
\end{solution}

\begin{solution}{pb:b2:deep-time-climate:1}
\textbf{1.} $1.3 + 1.3 + 4.5 + 4.2 = \qty{11.3}{Tmol.yr^{-1}}$, $\qty{0.136}{PgC.yr^{-1}}$.
\textbf{2.} $11.5\times0.012011 = \qty{0.138}{PgC.yr^{-1}}$.
\textbf{3.} Keduanya sesuai dalam batas sekitar \qty{2}{\percent}: pelapukan silikat mengimbangi pelepasan gas
vulkanik, seperti yang dituntut \omterm{def:b2:deep-time-climate:carbon-cycle}{siklus karbon jangka panjang}.
\textbf{4.} $\num{5.1e18}/0.02897 = \qty{1.76e20}{mol}$ udara; $\num{591e15}/12.011 =
\qty{4.92e16}{mol}$ karbon; $\num{4.92e16}/\num{1.76e20} = \num{2.80e-4}$, \qty{280}{ppm}.
\textbf{5.} $591/0.136 = \qty{4350}{yr}$.
\textbf{6.} $1 + 0.08 = 1.08$: hanya \qty{8}{\percent} lebih banyak.
\textbf{7.} $63000/(8.314\times288^2) = \qty{0.0914}{K^{-1}}$, yang dibulatkan menjadi 0.09.
\textbf{8.} $e^{0.09} = 1.09$: \qty{9}{\percent} lebih banyak per kelvin.
\textbf{9.} $\Delta T = S\log_2(1 + fM/A_0)$; turunannya pada 0 adalah $Sf/(A_0\ln2) =
3\times0.08/(591\times0.693) = \qty{5.86e-4}{K.PgC^{-1}}$.
\textbf{10.} $W_0k\times\num{5.86e-4} = 0.138\times0.09\times\num{5.86e-4} = \qty{7.28e-6}{yr^{-1}}$.
\textbf{11.} Kebalikannya adalah $A_0\ln2/(fkSW_0)$, yaitu $\tau$ pada proposisi itu, karena
persamaan yang dilinearkan adalah $\mathrm dM/\mathrm dt = -(\mathrm dW/\mathrm dM)\,M$.
\textbf{12.} Keempatnya: $\tau$ berbanding terbalik dengan masing-masing.
\textbf{13.} $1/\num{7.28e-6} = \qty{1.37e5}{yr}$.
\textbf{14.} $137\times0.09/0.119 = \qty{104}{ka}$.
\textbf{15.} $k_T = 0.107$: $137\times0.09/0.107 = \qty{115}{ka}$.
\textbf{16.} Lebih pendek dengan faktor 1.5 sampai 7: model satu kotak dan $f$-nya yang konstan membuat
termostat terlalu cepat.
\textbf{17.} $\num{50e6}/\num{1.37e5} = 364$: termostat sempat menyeimbangkan siklus itu pada setiap
tahap, sehingga penurunan itu mengikuti perubahan lambat pelepasan gas atau kemampuan lapuk.
\textbf{18.} $16\times0.137 = \qty{2.2}{ppm}$.
\textbf{19.} $0.08\times5000 = \qty{400}{PgC}$; $278\times991/591 = \qty{466}{ppm}$.
\textbf{20.} $3\log_2 1.677 = \qty{2.24}{K}$.
\textbf{21.} $0.138\,(e^{0.202} - 1) = \qty{0.0308}{PgC.yr^{-1}}$, \qty{22.3}{\percent}.
\textbf{22.} $137\times\ln10 = \qty{315}{ka}$.
\textbf{23.} Pelapukan mungkin lebih peka (limpasan, energi aktivasi yang lebih besar, umpan balik
lambat dari sensitivitas sistem Bumi yang lebih besar); dan karbon organik di darat dan di laut mungkin telah tumbuh
kembali dan mengambil karbon dari udara lebih cepat daripada pelapukan silikat.
\textbf{24.} Dalam model ini \omterm{def:b2:deep-time-climate:thermostat}{termostat pelapukan silikat} memindahkan karbon berlebih dengan
waktu e-folding sekitar \textbf{\qty[text-series-to-math=true]{137}{} ribu tahun}.
\end{solution}
